\begin{trapbox}
	\begin{description}[style=nextline, leftmargin=2.8em, labelsep=0.5em, itemsep=0.35em]
		\item[\textbf{\textcolor{CTrap}{易错点一（分母为零）：}}] 当 $\sum (x_i-\bar{x})^2 = 0$ 时，斜率公式的分母为零，无法计算 $\hat{b}$，但部分学生仍强行代入公式得到无意义结果。
		\item[\textbf{\textcolor{CTrap}{正确理解（x必须变异）：}}] 使用回归直线的前提是自变量 $x$ 的取值不完全相同，即一定有 $\sum (x_i-\bar{x})^2 > 0$；否则回归直线不存在（斜率未定义）。
		\item[\textbf{\textcolor{CTrap}{错因分析（忽略假设）：}}] 最小二乘法推导中假设分母非零，实际应用中若 $x$ 无变异则无法拟合，应检查数据质量。
	\end{description}

	\medskip
	\begin{description}[style=nextline, leftmargin=2.8em, labelsep=0.5em, itemsep=0.35em]
		\item[\textbf{\textcolor{CTrap}{易错点二（变量混淆）：}}] 计算 $\hat{b}$ 时错误地将分子写成 $\sum (y_i-\bar{y})^2$ 或分母写成 $\sum (y_i-\bar{y})^2$，导致斜率方向或大小错误。
		\item[\textbf{\textcolor{CTrap}{正确理解（回归方向）：}}] 回归分析中 $y$ 因变量、$x$ 自变量，$\hat{b}$ 公式分子是 $x$ 离差与 $y$ 离差的乘积之和，分母是 $x$ 离差平方和，不可对调。
		\item[\textbf{\textcolor{CTrap}{错因分析（概念不清）：}}] 混淆了相关系数与回归系数的定义，或是将 $x$ 与 $y$ 的地位等同。
	\end{description}

	\medskip
	\begin{description}[style=nextline, leftmargin=2.8em, labelsep=0.5em, itemsep=0.35em]
		\item[\textbf{\textcolor{CTrap}{易错点三（过点误解）：}}] 误以为回归直线一定经过某个原始数据点 $(x_i,y_i)$，从而用该点求回归方程。
		\item[\textbf{\textcolor{CTrap}{正确理解（过中心点）：}}] 回归直线唯一确定经过样本中心 $(\bar{x},\bar{y})$，但不一定经过任意观测点。
		\item[\textbf{\textcolor{CTrap}{错因分析（直观错觉）：}}] 由于直线拟合数据，学生直觉认为它应穿过数据点，但最小二乘法只保证整体误差最小而非经过某点。
	\end{description}
\end{trapbox}
